\pdfoutput=1
\documentclass[11pt]{article}
\usepackage{mathfin-paper}

\title{A Formally Verified Library of\\ Mathematical Finance in Lean~4}
\author{Raphael Coelho\thanks{Independent researcher. ORCID
  \href{https://orcid.org/0009-0001-6601-1023}{0009-0001-6601-1023}.
  Correspondence: \href{mailto:raphaelrrcoelho@gmail.com}{\texttt{raphaelrrcoelho@gmail.com}}.
  Artifact: \url{https://github.com/formal-applied-math/formal-mathfin}, release \texttt{v1.5.0}.}}
\date{October 2026}

\hypersetup{pdftitle={A Formally Verified Library of Mathematical Finance in Lean 4},
  pdfauthor={Raphael Coelho}}

\begin{document}
\maketitle

\begin{abstract}
\noindent
We describe a library of mathematical finance in the Lean~4 proof assistant, built on Mathlib and the
\texttt{BrownianMotion} package. Its benchmark records 535 results from the literature, each with a Lean
statement, a status, and a published account of what the formalization leaves out: 504 are proved in full,
18 restate an upstream theorem, and 13 are reduced cores that prove less than the result they are named
after. The library is about 100,000 lines of Lean, contains no \texttt{sorry}, and a build-time audit pins
every library constant that a benchmark proof cites to Lean's three standard axioms. It reaches from
continuous-time foundations, where the It\^o integral is constructed rather than assumed and carries It\^o's
formula without growth conditions, martingale representation, Girsanov's theorem for bounded predictable
drifts and Lipschitz stochastic differential equations, through no-arbitrage theory and derivative pricing to
quantitative risk management, fixed income, portfolio theory and actuarial mathematics. Two things make it more
than a catalogue. One is depth: the objects of continuous-time pricing are constructed rather than posited, from
Brownian motion and an It\^o integral built as an isometry. The other is an audit of faithfulness that makes the logical status of every claim checkable,
and whose review has caught what mechanical gates cannot see, including a convergence theorem whose hypothesis
could never be satisfied. Nearly all of the mathematics is classical, and what a formal base over it produces is
certified unification of known results. The library also hosts two contributed developments, by Robert Martin,
on the American put's exercise boundary and on Durrleman's condition for SVI smiles, that address problems its own survey
had listed as open. The contribution is methodological and infrastructural.
\end{abstract}

%======================================================================
\section{Introduction}

Interactive theorem provers have moved, over the last decade, from isolated landmark proofs to broad, reusable
libraries. Mathlib~\cite{Mathlib2020} now carries a substantial measure-theoretic probability development:
conditional expectation, martingales, the Radon--Nikodym theorem, Doob's theorems, Gaussian measures, and now the
definition of pre-Brownian motion. Degenne and collaborators built Brownian motion itself on top of it,
in the \texttt{BrownianMotion} package~\cite{DegenneEtAl2025,BrownianMotionRepo}. Mathematical finance sits directly on this foundation: arbitrage-free pricing,
the Black--Scholes model~\cite{BlackScholes1973}, the fundamental theorem of asset
pricing~\cite{HarrisonKreps1979}, risk measurement. Little of it had been formalized.

What existed was narrow and deliberate. Echenim and Peltier formalized the binomial model and its completeness
in Isabelle/HOL~\cite{EchenimPeltier2017}, and with Guiol the Cox--Ross--Rubinstein pricing of European
derivatives~\cite{EchenimGuiolPeltier2020}. Pusceddu and Bartoletti formalized constant-product automated market
makers in Lean~4~\cite{PuscedduBartoletti2024}. Nagy gave a machine-verified derivation from It\^o's lemma to
Black--Scholes in Lean~4~\cite{Nagy2026}, a focused vertical that, on its author's account, leaves the continuous
It\^o integral structurally verified rather than constructed. More recently, Ushakov and Berdinsky priced European
and barrier options through a Henstock--Kurzweil path integral in Lean~\cite{UshakovBerdinsky2026}. Each of these
is self-contained, and none offers a way to state \emph{how faithful} a formalization is to the mathematics it
claims.

This paper describes a library that is broad, deep in its continuous-time foundations, and audited for
faithfulness. It covers thirteen areas, from the It\^o integral to quantitative risk management, and the state
reported here is release \texttt{v1.5.0} of October 2026. Its It\^o calculus and its fundamental theorems of asset
pricing are treated in detail in two companion papers~\cite{CoelhoIto,CoelhoFTAP}; this one describes the
library as a whole, the principles that keep it coherent, and the audit that keeps it honest.

\paragraph{Changes from the third version}
The corpus has grown from 330 to 535 entries and the library has nearly doubled, from about 52,000 to about 100,000
lines. New since the third version are martingale representation and market completeness, It\^o's formula for It\^o
processes, the agreement of the integral with the \texttt{BrownianMotion} package's stochastic integral,
jump-diffusions, a quantitative-risk-management benchmark of 65 entries, and two contributed developments
(\cref{sec:provenance}); some results the library already contained, among them It\^o's formula under exponential
growth and stochastic differential equations, are reported here for the first time. The audit now pins every cited
constant rather than a curated list. Several statements of the third version are corrected here: the gloss of
the reduced-core tier (\cref{sec:audit}), the list of instances of the Garman normal form and the account of the
survival identity (\cref{sec:architecture}), the status of the binomial limit, whose stated hypothesis was
unsatisfiable at the time (\cref{sec:vacuity}), and the description of the integral $\int B\,\dd B$
(\cref{sec:depth}).

%======================================================================
\section{Contributions}\label{sec:contributions}

\begin{enumerate}[label=(\arabic*)]
\item \emph{Breadth on one foundation} (\cref{sec:breadth}). To our knowledge the broadest formalization of
  mathematical finance in any proof assistant: thirteen areas and 535 catalogued results on one set of
  definitions, with cross-module reuse rather than per-instrument re-derivation.
\item \emph{Depth in continuous time} (\cref{sec:depth}). The It\^o integral is constructed as an isometry, and
  It\^o's formula, martingale representation, Girsanov's theorem and stochastic differential equations are built
  on it; the risk-neutral measure is constructed by a change of measure rather than assumed, and the discounted price
  is identified as an It\^o integral with its integrand named.
\item \emph{A faithfulness audit} (\cref{sec:audit}), the most transferable contribution. Every result carries a
  tier that records how its Lean statement relates to the claim, and a published disclosure of what it leaves
  out. Build-enforced gates pin axioms, forbid unproved or definitional shortcuts in full entries, keep
  verification evidence fresh, and stop prose from naming an object a statement only asserts to exist.
  \Cref{sec:vacuity} reports what the gates missed and how it was found.
\item \emph{An honest map} (\cref{sec:evaluation,sec:limitations}). The reduced cores, the limitations of the
  headline results and of the apparatus itself, and the provenance of the code,
  including third-party ports and AI assistance (\cref{sec:provenance}).
\end{enumerate}

%======================================================================
\section{Related work and positioning}\label{sec:related}

Formalized finance has appeared mostly at formal-methods venues. \Cref{tab:related} compares the nearest
precedents along the two axes the prior art leaves open. Nagy's derivation~\cite{Nagy2026} is the closest, and
the library builds on it directly: three modules (\file{Foundations/DiscreteIto.lean}, \file{FTAPTwoState.lean},
\file{ItoLemma.lean}) adapt its discrete It\^o, two-state fundamental theorem and structural-drift framework, with
attribution in their headers. Where Nagy machine-checks one It\^o-to-Black--Scholes path and stops at the
continuous It\^o integral, this library constructs that integral and sets the same derivation inside a
comprehensive, audited development. The work of Echenim, Guiol and Peltier~\cite{EchenimPeltier2017,EchenimGuiolPeltier2020}
is the established precedent for formalized pricing in discrete time. The library re-implements the
constant-product invariants of Pusceddu and Bartoletti~\cite{PuscedduBartoletti2024} for breadth, adapting their
mathematical content with attribution, and claims nothing new in automated market makers. Ushakov and
Berdinsky~\cite{UshakovBerdinsky2026} take a different analytic route to pricing, through the Henstock--Kurzweil path
integral rather than a stochastic integral.

\begin{table}[tbp]
\centering\small
\begin{tabularx}{\linewidth}{@{}>{\raggedright\arraybackslash}p{0.24\linewidth}>{\raggedright\arraybackslash}X>{\raggedright\arraybackslash}X>{\raggedright\arraybackslash}X@{}}
\toprule
 & this work & Nagy (2026) & Echenim et al.\ (2017--20) \\
\midrule
prover & Lean~4 & Lean~4 & Isabelle/HOL \\
scope & 13 areas, 535 results & one derivation (It\^o to Black--Scholes) & binomial and CRR models, completeness \\
continuous It\^o integral & constructed & structurally verified & not applicable (discrete time) \\
risk-neutral measure & derived by a change of measure & derived in a two-state market & discrete-time model \\
faithfulness audit & four tiers, per-entry disclosure, axiom gates & not part of the work & not part of the work \\
\bottomrule
\end{tabularx}
\caption{This work against the nearest formalized-finance precedents.}
\label{tab:related}
\end{table}

Mathlib and the \texttt{BrownianMotion} package are the upstream foundation. The library consumes Mathlib's measure
theory, conditional expectation, martingales and Gaussian machinery, and the package's construction of Brownian
motion, its simple processes and predictable $\sigma$-algebra, local martingales and axiomatic stochastic integral;
\cref{sec:evaluation} reports what has been contributed back. On the benchmark side, StochBench~\cite{StochBench2026}
collects 450 graduate stochastic-process problems in Lean and, lacking infrastructure such as the It\^o integral,
states 336 of them with the required properties as hypotheses. The library's benchmark is of a different kind: it
records results that are proved, and its reduced-core tier makes exactly that hypothesis-taking visible where it
remains.

%======================================================================
\section{What the library proves}\label{sec:breadth}

\Cref{tab:areas,tab:areas2} list the thirteen areas with representative results. Each one-line statement describes the Lean
statement, not the textbook theorem it is named after; where they differ, the difference is given.

\begin{table}[tbp]
\centering\footnotesize
\renewcommand{\arraystretch}{1.1}
\begin{tabularx}{\linewidth}{@{}>{\raggedright\arraybackslash}p{0.15\linewidth}>{\raggedright\arraybackslash}X@{}}
\toprule
Area & Representative results \\
\midrule
Probability & Doob's $L^p$ maximal inequality for non-negative submartingales (\lean{maximal_ineq_Lp}); convergence of
  right-continuous $L^p$-bounded martingales, $p>1$, almost surely along the integers and in measure along the reals;
  the Ionescu-Tulcea path measure of a Markov chain on cylinders. \\
Stochastic calculus & The It\^o integral as an isometry and as a continuous local martingale; It\^o's formula in $L^2$ under
  bounded or exponentially growing derivatives, without growth conditions as a local martingale
  (\lean{ito_formula_unrestricted}), and for It\^o processes with bounded adapted continuous coefficients
  (\lean{ito_formula_adapted}); martingale representation; Girsanov's theorem for bounded predictable drifts, at the
  level of finite-dimensional laws; Lipschitz SDEs on short horizons; the heat-equation link of Feynman--Kac~\cite{CoelhoIto}. \\
No-arbitrage & The fundamental theorem over several periods on a finite probability space with one risky asset
  (\lean{ftap_discrete}), and in one period with finitely many assets, on a finite state space and on any probability space
  (\lean{ftap_one_period_vector})~\cite{CoelhoFTAP}; an equivalent
  martingale measure excludes arbitrage by simple strategies in continuous time; change of num\'eraire; completeness of
  the Brownian market with a unique hedge (\lean{exists_replicating_strategy}); uniqueness of the pricing measure under a
  square-integrable density. \\
Option pricing & The Black--Scholes call under a lognormal hypothesis that is itself derived from the physical model; the
  Black--Scholes equation from Feynman--Kac (\lean{bsV_satisfies_bs_pde_via_feynmanKac}); Breeden--Litzenberger
  in second-order form for Black--Scholes and for jump-diffusions with a Gaussian part, and in first-order form (call prices
  at all strikes determine the law) for any law with a finite forward; jump-diffusions, assumed as a hypothesis structure:
  the discounted price is a martingale exactly at the compensated drift
  (\lean{JumpDiffusionProcess.martingale_iff}), Merton's 1976 series at every date before maturity
  (\lean{condExp_call_eq_mertonCallPrice}), Esscher pricing, and incompleteness at one date
  (\lean{exists_call_esscher_ne_merton}). \\
Exotic options & The $n$-date geometric Asian call in Black--Scholes form; Margrabe's formula as a Black--Scholes call on the
  price ratio when the ratio is lognormal, and the change-of-num\'eraire identity for the exchange payoff; quanto forwards; variance swaps, with realized variance converging to $\sigma^2T$ in $L^2$ for any drift. \\
Lattice models & The Cox--Ross--Rubinstein call price converges to the Black--Scholes price
  (\lean{binomialPrice_call_tendsto_bs_closed}), with a lemma that its no-arbitrage hypothesis holds when
  $\abs r\sqrt T<\sigma$ (\lean{binomialNoArb_crr}); American options as Snell envelopes; Andr\'e's reflection principle as a
  bijection of lattice paths. \\
\bottomrule
\end{tabularx}
\caption{Representative results in probability, stochastic calculus and pricing. Every name is a Lean declaration in
the library.}
\label{tab:areas}
\end{table}

\begin{table}[!t]
\centering\footnotesize
\renewcommand{\arraystretch}{1.1}
\begin{tabularx}{\linewidth}{@{}>{\raggedright\arraybackslash}p{0.15\linewidth}>{\raggedright\arraybackslash}X@{}}
\toprule
Area & Representative results \\
\midrule
Fixed income, credit & The Vasicek bond price in affine form, with the integrated rate defined through its Wiener
  representation; the forward measure as a change of num\'eraire, in a constant-rate model where it coincides with the
  risk-neutral measure; duration and convexity as price derivatives; hazard rates as logarithmic derivatives of survival;
  the CDS fair spread from an assumed balance of premium and protection legs; under independence, the first-to-default
  hazard as the sum of the single-name hazards; the KMV--Merton default probability through the distance to default. \\
Portfolio theory & CAPM from the equilibrium first-order condition; the Black--Litterman posterior mean as the unique solution of its
  precision-weighted normal equation, without the Bayesian derivation; Markowitz critical points
  from the Lagrangian; Kelly growth over $n$ periods; the Kelly portfolio as num\'eraire induces the risk-neutral measure in a
  two-outcome market. \\
Risk measures & The Artzner--Delbaen--Eber--Heath representation of coherent risk measures on a finite state space
  (\lean{coherentRisk_isLUB}); value-at-risk as the quantile of an arbitrary law; the Rockafellar--Uryasev formula, coherence
  and dual representation of expected shortfall for every integrable loss; elicitability of value-at-risk, on laws with a
  finite mean and a unique quantile, and its failure for expected shortfall. \\
Dependence, extremes & Sklar's theorem for continuous margins; Fr\'echet--Hoeffding bounds; max-stability of the GEV law;
  the generalized Pareto law in its domain of attraction; maxima of Poisson-many losses. Fisher--Tippett and
  Pickands--Balkema--de~Haan are not formalized. \\
Credit portfolios & The value-at-risk of the one-factor Gaussian default fraction converges to the asymptotic
  single-risk-factor quantile, the worst-case default rate at the core of the Basel IRB formula
  (\lean{tendsto_valueAtRisk_oneFactorLossFraction}). \\
Microstructure & A positive Glosten--Milgrom spread from trader-mix primitives; for one asset, the quadratic ansatz of
  Bergault--Evangelista--Gu\'eant--Vieira solves the approximate (linear--quadratic) Hamilton--Jacobi equation whenever its
  coefficients solve the associated ordinary differential equations, and for several assets the matrix Riccati equation of
  the leading coefficient is verified; the true value function is out of scope. \\
Other & A contract language that separates payoffs from pricing models, after Bilokon~\cite{Bilokon2026}; survival models,
  annuity-certain closed forms and compound Poisson losses; constant-product AMMs; strict convexity of the American put's
  exercise boundary and a polynomial-sign characterisation of Durrleman's condition for raw SVI, both by Robert Martin, the
  latter extended by the library to butterfly-free smiles (\cref{sec:provenance}). \\
\bottomrule
\end{tabularx}
\caption{Representative results in the applied areas. Every name is a Lean declaration in the library.}
\label{tab:areas2}
\end{table}

No individual theorem in the table is new; the mathematics is classical throughout, with the exception of the two
contributed developments discussed in \cref{sec:provenance}. The contribution is the development as a whole: one
foundation under all of it, depth where earlier efforts assume, and an audited statement of what each result proves.

%======================================================================
\section{Depth: the continuous-time spine}\label{sec:depth}

The library's depth comes from carrying stochastic integration through to a genuine It\^o integral, the foundation
earlier formalizations assume or sketch, and then consuming it. The companion paper~\cite{CoelhoIto} develops this
in full; here is the spine.

\paragraph{The integral}
The It\^o integral $\ito_T$ is a linear isometry from the square-integrable predictable integrands on $(0,T]\times\Omega$
into $L^2(\mu)$ (\lean{itoIntegralCLM_T}), the unique continuous extension of the elementary integral on simple adapted
processes, obtained by a $\pi$--$\lambda$ density argument and Mathlib's \lean{LinearMap.extendOfNorm}. As a process it is
the conditional projection of its terminal value, hence a continuous $L^2$ martingale, and it has a continuous
modification that is a local martingale on the half-line. Its range together with the constants is all of
$L^2(\F^B_T)$, which is martingale representation, and it satisfies the \texttt{BrownianMotion} package's axiomatic
characterisation of a stochastic integral. One worked example is the identity behind $\int_0^TB\,\dd B$: an
integrand $g_B$, constructed as the limit of truncated step processes of $B$, has $\ito_T(g_B)=\frac12(B_T^2-B_0^2-T)$
(\lean{itoIntegralCLM_T_brownian}). That statement does not itself name $g_B$ as $B$; the named form is an instance
of It\^o's formula under exponential growth, which does name its integrand, though that instance is not stated
separately~\cite{CoelhoIto}.

\paragraph{It\^o's formula}
In $L^2$, for $f(t,x)$ whose partial derivatives up to $f_{tt}$, $f_{tx}$ and $f_{xxx}$ are bounded or grow at most like
$e^{\lambda\abs x}$, $f(T,B_T)-f(0,B_0)=\ito_T(f_x(\cdot,B))+\int_0^T(f_t+\frac12f_{xx})(s,B_s)\,\dd s$ with the integrand
named. For every $C^3$ function, with no growth condition, the compensated process is a continuous local martingale on
the null-augmented filtration, by localization at exit times (\lean{ito_formula_unrestricted}). For an It\^o process
$X=x_0+\int b\,\dd s+\int\sigma\,\dd B$ with bounded adapted continuous coefficients and $f$ with bounded $f'$ and $f''$, the
formula holds at a fixed time (\lean{ito_formula_adapted}).

\paragraph{The risk-neutral measure, derived}
In most treatments, and in earlier formalized pricing, an equivalent martingale measure is posited and prices are its
expectations. Here it is constructed. A static Girsanov change of measure by an Esscher density turns the physical
Gaussian driver into one under which the terminal asset has exactly the lognormal law the Black--Scholes layer takes as
its hypothesis, so that hypothesis is a theorem (\lean{BSCallHyp.exists_of_physical}) and the call formula holds under a
measure constructed from the physical Gaussian driver alone (\lean{bs_call_formula_of_physical}). The dynamic Girsanov theorem is proved along a ladder of
drifts (constant, simple adapted, continuous adapted, bounded predictable): under
$\frac{\dd Q}{\dd\mu}=\exp(-\ito_T(\theta)-\frac12\int_0^T\theta^2\,\dd s)$ the process $B+\int\theta\,\dd s$ starts at zero and
has independent $N(0,t-s)$ increments on $[0,T]$, which fixes its finite-dimensional laws; path continuity and
independence from the past are not stated. The discounted geometric Brownian motion is an It\^o integral with its integrand named
(\lean{discountedGBM_eq_itoIntegral}), so the pricing layer meets the continuous world through the integral itself,
not only through the marginal laws that the third version described.

\paragraph{Completeness and its limits}
Every square-integrable $\F^B_T$-measurable claim is $\E[H]+\int_0^T\varphi\,\dd B$ for a unique hedge
(\lean{exists_replicating_strategy}); integrating against the price $S=S_0+\int\sigma\,\dd B$ instead, with $\sigma\ne0$
almost everywhere, the hedge is a holding in the price. A probability measure with square-integrable density under
which $S$ is a martingale agrees with $\mu$ on $\F^B_T$: the direction of the second fundamental theorem from
completeness to uniqueness. The converse, the general Girsanov theorem under Novikov's condition, and the identification
of the hedge (Clark--Ocone) are open.

\paragraph{Equations and jumps}
Lipschitz stochastic differential equations have a unique solution in the predictable $L^2$ space whenever
$TL_b+\sqrt TL_\sigma<1$ (\lean{picardMap_exists_unique_fixedPoint}), realized as a sample-path process
(\lean{sde_pathwise_decomposition}). Jump-diffusions
enter through a hypothesis structure (\lean{JumpDiffusionProcess}): an adapted process with independent increments whose
laws are those of drift plus Gaussian plus compound Poisson. On it the discounted price is a martingale exactly at the
compensated drift, the conditional call price at every date is Merton's series~\cite{Merton1976}, and two equivalent
compensated laws at one date price some call differently. Such a process is constructed only without jumps, as Brownian
motion with drift.

\paragraph{From lattice to continuum}
The Cox--Ross--Rubinstein call price converges to the Black--Scholes price, by characteristic functions and L\'evy's
continuity theorem, with put--call parity removing the uniform-integrability obstruction. The limit is proved for call
prices and for the law of the terminal log-return; Donsker's invariance principle is not formalized.

%======================================================================
\section{Architecture: principles and bridges}\label{sec:architecture}

The library stays coherent by deriving results from a few reusable principles rather than re-proving each instrument.
Four pillars carry the mathematics, and a small set of bridges connects them.

\paragraph{I. No-arbitrage as convex duality}
The finite fundamental theorem is a separating hyperplane between the attainable gains and the simplex; the
Artzner--Delbaen--Eber--Heath representation of coherent risk measures~\cite{ADEH1999} separates a rejected position from
the acceptance cone. The two separations are sibling lemmas in one module, proved from the same two helper
lemmas~\cite{CoelhoFTAP}. On a finite state space, pricing as a non-negative linear functional on payoffs gives put--call
parity, the forward price and convexity of the call in the strike as corollaries of one structure rather than
instrument-by-instrument results.

\paragraph{II. Stochastic calculus}
Continuous-time models are built on Brownian motion, and the It\^o integral and formula of \cref{sec:depth} are the
calculus for their functionals; pricing hypotheses stated on marginal laws are tied to Brownian motion by the bridges
below.

\paragraph{III. The probabilistic--analytic duality}
A price is both a risk-neutral expectation and the solution of a partial differential equation, and Feynman--Kac is the
bridge. The Black--Scholes equation is derived by writing the price as a heat-kernel integral and using the heat equation
for the kernel, without differentiating the closed form.

\paragraph{IV. Exponential families and intensities}
Closed forms come from Gaussian, lognormal and exponential structure. One Esscher module serves three users: its tilt,
Mathlib's \lean{Measure.tilted} at a linear exponent, gives the static Girsanov change of measure and Esscher pricing of
jump-diffusions, and its moment-generating-function uniqueness lemma identifies Brownian increments under a new measure.
And one exponential of an integrated intensity, \lean{survivalFromIntensity}, is the shared definition of the
intensity-based actuarial and credit survival functions.

\paragraph{The Garman normal form}
Many Black--Scholes-family prices have the form $A\,\Phi(d_1)-K\cdot\mathrm{DF}\cdot\Phi(d_2)$. The library proves this
equality for the standard call, Black-76 futures options, payer swaptions, the call with a continuous dividend yield, and
Margrabe's exchange option~\cite{Margrabe1978}. The third version listed the put,
digital and power options among the proved instances; they are not. Garman--Kohlhagen is an instance only through the
dividend-yield case, without a theorem of its own.

\paragraph{Grounding in Brownian motion}
Pricing hypotheses are stated marginally (\lean{BSCallHyp} asks only that the terminal asset be lognormal under the pricing
measure), and bridge modules show that they follow from a pre-Brownian motion, through which the flagship prices connect
to the continuous-time object.

\paragraph{Certified unification}
What a formal base over classical mathematics produces is not new theory but machine-checked identities between objects
developed for different purposes. For Gaussian portfolios with a common mean, minimising value-at-risk or expected shortfall
at a level above one half is equivalent to minimising variance (\lean{isMinOn_valueAtRisk_iff_isMinOn_portfolioVarN}). The
$N$-asset portfolio variance with uncorrelated assets of common variance is that variance times the Herfindahl concentration
index (\lean{portfolioVarN_diag_eq_herfindahl}), and with equal correlation it is an explicit mixture of the index and the
squared total weight. Mortality and reduced-form credit share one intensity-based survival function. The third version presented this last identity
as one between independently developed definitions; in fact both are defined through a shared
\lean{survivalFromIntensity}, so it holds by definition, and it is the shared definition, not the identity, that records the
unification.

\paragraph{A representation layer}
The contract layer (\file{Contracts/}) is not a fifth principle but a representation across the four: payoffs become
inductive objects, proved measurable and adapted to their own observation times, and for the European call, put, digital and
capped call under single-asset Black--Scholes their values reduce to the closed forms the pricing layer already proves. Its design follows Bilokon's \emph{The Contract Is Not the Model}~\cite{Bilokon2026}; no code
is taken from it.

%======================================================================
\section{The faithfulness audit}\label{sec:audit}

Formalization papers often leave implicit the gap between ``we proved $X$'' and ``we proved a weaker $X$ under added
hypotheses''. The library makes the gap explicit and partly machine-checkable.

\subsection{Tiers and disclosures}
Every benchmark entry declares one of four statuses. \status{full}: the Lean statement is the mathematical claim, derived
from foundational primitives. \status{library\_wrapper}: a thin restatement of an upstream Mathlib or \texttt{BrownianMotion}
theorem. \status{reduced\_core}: an honest but narrower statement; in twelve of the thirteen current cases the textbook
conclusion is a field of a structure and the theorem reads it off by projection, so nothing is derived, and in the
thirteenth a special case is proved. \status{placeholder}: none remain. Delivery claims count only \status{full} and
\status{library\_wrapper}. (The third version glossed the reduced core as ``holds under an added hypothesis or with an
axiomatized sub-step'', which understated how little the twelve specifications prove.) Each entry's \texttt{description}
states what the entry proves, not the textbook theorem it is named after, and its \texttt{formalization\_scope} gives the
long-form disclosure of what it leaves out. Both are published, together with the Lean statement, in a public
dataset~\cite{HFDataset}.

\subsection{Machine-checked gates}
The tiers are metadata, checked by test gates that continuous integration runs before the Lean build, and the Lean build
itself checks axioms.
\begin{itemize}
\item \emph{Axiom audits.} A curated file pins the \texttt{\#print axioms} output of 372 headline constants under
  \texttt{\#guard\_msgs}. A generated file pins 600: every library constant cited in proof position by any benchmark
  snippet (571), resolved by declaration through \texttt{open}s, dot notation and foreign namespaces, and the 29 upstream
  constants that wrapper entries re-export. Together they pin 760 distinct constants. The upstream pins matter because \texttt{BrownianMotion} contains
  \texttt{sorry}s at the pin; none lies beneath a pinned constant. Any axiom beyond \texttt{propext},
  \texttt{Classical.choice} and \texttt{Quot.sound}, any \texttt{sorryAx}, or a renamed constant fails the build.
\item \emph{Text gates.} After comment stripping, library sources contain no \texttt{sorry}, \texttt{admit},
  \texttt{native\_decide}, \texttt{polyrith}, search-tactic suggestions or premise-selection calls.
\item \emph{Shape gates.} A full entry may not bundle an assumed \texttt{Prop}-valued structure (the reduced-core pattern),
  nor rest on a definitional \texttt{rfl} proof; every reduced-core description must say what is not derived; and an
  existential statement whose witness is never pinned down may not be described by prose that names an integral of it.
\item \emph{Evidence freshness.} A verification ledger records, for every entry, a hash of its snippet, the transitive library
  modules it imports and the toolchain pins under which it last verified; an entry whose inputs change becomes stale and must
  re-verify.
\end{itemize}

\subsection{What the gates cannot see}\label{sec:vacuity}
Two failure modes pass every mechanical gate. The first is vacuity. The Cox--Ross--Rubinstein convergence theorems
that the third version of this paper reported as checked end to end carried the hypothesis that the $n$-step tree is
arbitrage-free for every $n$,
which fails at $n=0$, where the step size is $T/0=0$; so the theorems were true of no market and said nothing. No axiom
audit, kernel replay or \texttt{sorry} gate can detect this. It was found by reading the statement, not by any gate,
and fixed by quantifying over $n\ge1$ and proving that the hypothesis then holds whenever
$\abs r\sqrt T<\sigma$ (\lean{binomialNoArb_crr}). The external check of \cref{sec:external} likewise pairs its
theorem with a satisfiability witness.

The second is prose that outruns its statement: the build is green precisely because the statement is weaker than the
description. The library has repeatedly found such drift in its own text, including the Girsanov results, once described as
producing a Brownian motion when they fix finite-dimensional laws, and the integral $\int B\,\dd B$ of \cref{sec:depth}. Only
its mechanical slice is gated. The rest is a review discipline: every change to proof content closes with a multi-agent
review against eight lenses, whose standing first pass reads every touched description, docstring and paragraph against its
statement, and whose cadence (one review per twelve new entries at most) is itself test-enforced. This paper was revised with
the same first pass; the corrections listed in the introduction are what it found.

\subsection{External checks}\label{sec:external}
For the coherent-risk representation theorem the repository carries a Mathlib-only restatement, with a witness that coherent
risk measures exist, and a proof of both from the library that every build re-checks; the external checker
Comparator~\cite{Comparator} is configured to confirm that the proved statements are exactly the restated ones and that only the
three standard axioms are used. A kernel replay of every declaration through \texttt{leanchecker} exists as a workflow but has never completed: the full Mathlib
environment it loads does not fit in the memory of a hosted runner.

%======================================================================
\section{Evaluation}\label{sec:evaluation}

\begin{table}[tbp]
\centering\small
\begin{tabular}{@{}lr@{}}
\toprule
Benchmark entries & 535 \\
\quad \status{full} & 504 \\
\quad \status{library\_wrapper} & 18 \\
\quad \status{reduced\_core} & 13 \\
\quad \status{placeholder} & 0 \\
Delivery-ready (\status{full} $+$ \status{library\_wrapper}) & 522 \\
\midrule
Lean files / lines under \file{MathFin/} (excluding audit and blueprint files) & 560 / 100{,}315 \\
\quad of which contributed ports (\cref{sec:provenance}) & 214 / 24{,}238 \\
Theorem and lemma declarations & 3{,}841 \\
Constants pinned by the axiom audits (curated / generated / distinct) & 372 / 600 / 760 \\
Verification-ledger rows (all fresh) & 535 \\
\bottomrule
\end{tabular}
\caption{The library at release \texttt{v1.5.0}.}
\label{tab:numbers}
\end{table}

\Cref{tab:numbers} gives the numbers, from the repository's coverage report and source tree. The reduced cores sit at the
genuinely hard frontier of continuous-time and discrete-process theory: the reflection principle, nowhere differentiability and
the law of the iterated logarithm for Brownian motion; Novikov's condition, the general Girsanov theorem, L\'evy's characterisation
and the two-dimensional It\^o formula; five theorems on Markov chains (the strong Markov property, recurrence criteria, uniqueness
of the stationary distribution, the ergodic theorem and convergence to equilibrium); and, as a proved special case, the exponential
law of the first interarrival time of a Poisson process. The mathematical-finance benchmark, the largest, has 336 of its 337 entries
full. The library is most faithful where the mathematics is classical and elementary, and its partial cores sit where the
infrastructure is missing.

\paragraph{Reproducibility}
The toolchain is pinned: Lean \texttt{v4.33.0-rc1}~\cite{deMouraUllrich2021}, Mathlib at \texttt{0434c03} and
\texttt{BrownianMotion} at \texttt{0d5b6eb}, in \file{lean-toolchain} and \file{lake-manifest.json}. A plain \texttt{lake build}
from the repository root checks every proof; release \texttt{v1.5.0} is archived on Zenodo~\cite{Zenodo}, and a pinned container
image (\texttt{ghcr.io/formal-applied-math/mathfin-verify}) reproduces the build.

%======================================================================
\section{Provenance}\label{sec:provenance}

A library that audits its claims should also account for its code. The repository records provenance file by file and in a
machine-readable self-report.

\paragraph{Contributed developments}
About a quarter of the source lines are ports of two developments by Robert Martin, each kept under its original authorship. The first
proves, in the constant-parameter Black--Scholes model with $r>0$ and dividend yield $0\le q\le r$, that the American put's
early-exercise boundary $B(\tau)$ is strictly convex in the time to expiry and that $\log(B(\tau)/K)$ is
convex~\cite{MartinAmericanPut}. The value is the supremum over all stopping times of the completed Brownian filtration, and the
headline theorems assume no classical solution of the free-boundary problem; a further, conditional part of the port, stated over
an assumed classical solution that nothing shows inhabited, is not claimed. The second gives a fixed Boolean combination of signs
of polynomials in the five raw-SVI parameters that holds exactly when the parameters are admissible and Durrleman's condition
holds, that is, when the call priced through the SVI smile is convex in the strike~\cite{MartinSVI,GatheralJacquier2014}. The
library adds the large-strike tail condition: with one further sign the formula holds exactly when that call is a normalised call
price function in Roper's sense, which is freedom from butterfly arbitrage~\cite{Roper2010}; recovering the risk-neutral law from
such a function is not formalized.
The repository's survey of open problems had listed both questions as open, the first for $0<q<r$ and the second as an explicit
semialgebraic description of the SVI domain, and records both as resolved by these developments. The credit is Robert Martin's;
the first was, on its author's report, substantially AI-assisted; neither has been refereed; and their status is that of
machine-checked, axiom-pinned results.

\paragraph{Other sources}
Nagy's derivation~\cite{Nagy2026} and the automated-market-maker formalization of Pusceddu and Bartoletti~\cite{PuscedduBartoletti2024}
are adapted in four modules with attribution headers. The contract layer takes design ideas, but no code, from
Bilokon~\cite{Bilokon2026}. Yosuke Ito's Isabelle entry on actuarial mathematics was consulted, with permission, for the survival
layer. Alfredo Garcia contributed the Glosten--Milgrom model and the drifted price path, and Aleksey Salkutsan downside performance
metrics and forward-rate agreements. Much of the benchmark follows Saporito's lecture notes on stochastic processes~\cite{Saporito},
the corpus spine.

\paragraph{AI assistance}
The library was written with AI assistance, and the repository says so. The author used an AI coding assistant (Claude Code) as an
authoring assistant throughout. Four benchmark entries came from an automated pipeline in which language models drafted and proved the
statements; each was reviewed and refined by hand, and the pipeline is treated as a scout, not an author. The gates of
\cref{sec:audit} apply to all code unchanged, whoever or whatever wrote it, and every proof is checked by the Lean kernel when the
library builds.
This manuscript was revised with the same assistant, and its statements were checked against the Lean sources of release
\texttt{v1.5.0}.

%======================================================================
\section{Limitations}\label{sec:limitations}

The limitations of the headline results are those their statements carry. Girsanov's theorem is proved for bounded predictable
drifts and at the level of finite-dimensional laws. Of the second fundamental theorem only the direction from completeness to
uniqueness is proved, and martingale representation proves that a hedge exists and is unique without identifying it. It\^o's formula
for an It\^o process is proved at a fixed time for bounded adapted continuous coefficients and $C^2$ functions $f$ with bounded first
and second derivatives. Stochastic differential equations are solved only on horizons with $TL_b+\sqrt TL_\sigma<1$. The jump-diffusion results
assume the process (\lean{JumpDiffusionProcess}). The multi-period fundamental theorem, for one risky asset, and the representation of
coherent risk measures are proved on finite probability spaces; the general multi-period theorem of Dalang, Morton and Willinger is open. Several
results are narrower than their names: the Margrabe grounding asserts the existence of its pricing measure and driver, the
Vasicek integrated rate is defined through its Wiener representation, and the market-making result solves only the approximate
Hamilton--Jacobi equation.

The apparatus has limits of its own. Library theorems that no benchmark entry cites and the curated audit omits are not axiom-pinned;
the text gate and the Lean kernel still apply to them. Continuous integration checks that the ledger is fresh, not that it was
re-verified; re-verification runs on demand. The public dataset publishes descriptions and scope disclosures but not the provenance
notes. And the kernel replay has never completed.

%======================================================================
\section{Conclusion}\label{sec:conclusion}

The library is, to our knowledge, the most comprehensive machine-checked development of mathematical finance to date, but its real
contributions are not its size. They are its depth, a constructed It\^o integral on which the objects of continuous-time pricing are
derived, and an audit that makes every claim's logical status checkable and keeps prose from outrunning statements.

What the library is \emph{not} is equally part of the point. Machine-checking classical financial mathematics yields certified
unification of known results, not new financial theory, and the tiers, disclosures and gates make that honesty checkable. A formal
base can nonetheless host new results once they exist: two contributed developments address questions its own survey had listed as
open, and the library's apparatus is what lets a reader see exactly what they prove and on what they rest. The open frontier is the
natural next layer for a community library: the general multi-period Dalang--Morton--Willinger theorem and its continuous-time
counterpart, Girsanov's theorem under Novikov's condition, It\^o's formula for It\^o processes without bounds, pathwise quadratic
variation, the Clark--Ocone formula, and the thirteen results still classified as reduced cores.

%======================================================================
\begin{thebibliography}{99}
\small
\bibitem{ADEH1999} P.~Artzner, F.~Delbaen, J.-M.~Eber, D.~Heath. Coherent measures of risk. \emph{Mathematical Finance},
  9(3):203--228, 1999.
\bibitem{Bilokon2026} P.~Bilokon. The contract is not the model: proof-carrying exotic derivatives and the economics of model risk.
  Working paper, Imperial College London, 2026.
\bibitem{BlackScholes1973} F.~Black, M.~Scholes. The pricing of options and corporate liabilities. \emph{Journal of Political
  Economy}, 81(3):637--654, 1973.
\bibitem{BrownianMotionRepo} R.~Degenne et al. \emph{brownian-motion}: a Lean~4 development of Brownian motion and stochastic
  integrals. \url{https://github.com/RemyDegenne/brownian-motion}, 2024--2026.
\bibitem{CoelhoFTAP} R.~Coelho. The fundamental theorem of asset pricing, formalized in Lean~4. arXiv:2606.28990, 2026.
\bibitem{CoelhoIto} R.~Coelho. A machine-checked It\^o calculus for Brownian motion. arXiv:2606.15089, 2026.
\bibitem{Comparator} \emph{Comparator}: a checker for Lean proofs against Mathlib-only statements.
  \url{https://github.com/leanprover/comparator}, 2026.
\bibitem{DegenneEtAl2025} R.~Degenne, D.~Ledvinka, E.~Marion, P.~Pfaffelhuber. Formalization of Brownian motion in Lean.
  arXiv:2511.20118, 2025.
\bibitem{deMouraUllrich2021} L.~de~Moura, S.~Ullrich. The Lean~4 theorem prover and programming language. In \emph{CADE~28},
  LNCS 12699, pp.~625--635, 2021.
\bibitem{EchenimPeltier2017} M.~Echenim, N.~Peltier. The binomial pricing model in finance: a formalization in Isabelle. In
  \emph{CADE~26}, LNCS 10395, 2017. Extended version: \emph{Journal of Automated Reasoning}, 63, 2019.
\bibitem{EchenimGuiolPeltier2020} M.~Echenim, H.~Guiol, N.~Peltier. Formalizing the Cox--Ross--Rubinstein pricing of European
  derivatives in Isabelle/HOL. \emph{Journal of Automated Reasoning}, 64:737--765, 2020.
\bibitem{GatheralJacquier2014} J.~Gatheral, A.~Jacquier. Arbitrage-free SVI volatility surfaces. \emph{Quantitative Finance},
  14(1):59--71, 2014.
\bibitem{HarrisonKreps1979} J.~M.~Harrison, D.~M.~Kreps. Martingales and arbitrage in multiperiod securities markets.
  \emph{Journal of Economic Theory}, 20(3):381--408, 1979.
\bibitem{HFDataset} R.~Coelho. \emph{formal-mathfin-theorems} (dataset).
  \url{https://huggingface.co/datasets/formal-applied-math/formal-mathfin-theorems}, 2026.
\bibitem{Margrabe1978} W.~Margrabe. The value of an option to exchange one asset for another. \emph{Journal of Finance},
  33(1):177--186, 1978.
\bibitem{MartinAmericanPut} R.~Martin. \emph{AmericanPutConvexity}: a Lean development, with the manuscript ``Log-convexity of the
  exercise boundary of an American put option''. \url{https://github.com/robertmartin8/AmericanPutConvexity}, 2026.
\bibitem{MartinSVI} R.~Martin. \emph{ButterflyFreeSVI}: a Lean development. \url{https://github.com/robertmartin8/ButterflyFreeSVI},
  2026.
\bibitem{Mathlib2020} The mathlib Community. The Lean mathematical library. In \emph{CPP 2020}, pp.~367--381, 2020.
\bibitem{Merton1976} R.~C.~Merton. Option pricing when underlying stock returns are discontinuous. \emph{Journal of Financial
  Economics}, 3(1--2):125--144, 1976.
\bibitem{Nagy2026} T.~Nagy. From It\^o to Black--Scholes: a machine-verified derivation in Lean~4. SSRN Working Paper 6336503, 2026.
\bibitem{PuscedduBartoletti2024} D.~Pusceddu, M.~Bartoletti. Formalizing automated market makers in the Lean~4 theorem prover. In
  \emph{FMBC 2024}, OASIcs 118. Schloss Dagstuhl, 2024.
\bibitem{Roper2010} M.~Roper. Arbitrage free implied volatility surfaces. Preprint, University of Sydney, 2010.
\bibitem{Saporito} Y.~F.~Saporito. \emph{Stochastic Processes: Discrete and Continuous Time}. Lecture notes.
\bibitem{StochBench2026} I.~Davidovich, D.~Ganguly, V.~Singh, V.~Chaudhary. StochBench: a domain-specific benchmark for stochastic
  processes in Lean. arXiv:2609.09264, 2026.
\bibitem{UshakovBerdinsky2026} A.~S.~Ushakov, Y.~N.~Berdinsky. Henstock--Kurzweil path integral in financial mathematics: a
  machine-verified pricing of European and barrier options. arXiv:2608.19223, 2026.
\bibitem{Zenodo} R.~Coelho. \emph{formal-mathfin: a formally verified library of mathematical finance in Lean~4} (software). Zenodo,
  \url{https://doi.org/10.5281/zenodo.20477781}, 2026.
\end{thebibliography}

\end{document}
